% Notes (epub:type="note")

% This file defines the LaTeX rendering for semantic notes:
%   note
%   note.idea
%   note.insight       (minor insight)
%   note.insight.major (major insight)
%   note.warning
%   note.breadcrumb
%   note.contrast
%   note.navigation
%   note.bridge
%   note.misuse
%   note.diagnostic
%   note.closing
%
% It is built around a small presentation framework:
%   \noteheaderinline{glyph}        – inline glyph before first sentence
%   \noteheaderornament{glyph}      – standalone ornamental row (e.g., ⁂)
%   \noteheaderbadgenormal{glyph}   – badge strip with neutral note color
%   \noteheaderbadgewarning{glyph}  – badge strip with warning color
%
% The Lua filter inserts:
%   \notevariantideainline
%   \notevariantinsight
%   \notevariantinsightmajor
%   \notevariant{<variant>}
%
% which are mapped here onto the presentation primitives.

%--------------------------------------------------------------------
% Glyphs
%--------------------------------------------------------------------
% Glyph font assumed defined somewhere, e.g.:
%   \newfontfamily\noteglyphfont{Noto Emoji}
% or similar.

\newcommand{\glyphnote}{{\noteglyphfont ∴}}      % neutral/default note
\newcommand{\glyphidea}{{\noteglyphfont 💡}}     % note.idea
\newcommand{\glyphinsight}{{\noteglyphfont ⁂}}  % note.insight
\newcommand{\glyphwarning}{{\noteglyphfont ※}}  % note.warning (calmer than ⚠)
\newcommand{\glyphbreadcrumb}{{\noteglyphfont ↗}} % note.breadcrumb (navigational cue)
\newcommand{\glyphcontrast}{{\noteglyphfont ⇆}}  % note.contrast
\newcommand{\glyphnavigation}{{\noteglyphfont ⇢}} % note.navigation
\newcommand{\glyphbridge}{{\noteglyphfont ⟷}}   % note.bridge
\newcommand{\glyphmisuse}{{\noteglyphfont ✕}}    % note.misuse
\newcommand{\glyphdiagnostic}{{\noteglyphfont 🔍}} % note.diagnostic
\newcommand{\glyphclosing}{{\noteglyphfont ∎}}   % note.closing


%--------------------------------------------------------------------
% Presentation primitives (namespaced as \noteheader...)
%--------------------------------------------------------------------

% 1. Inline header: small glyph before the first sentence.
\newcommand{\noteheaderinline}[1]{%
  \raisebox{0.05em}{#1}\hspace{0.4em}\ignorespaces
}

% 2. Ornamental header: a standalone glyph row.
%    Using your "lightly indented asterism" style.
% Minor insight (note.insight): small ornament row
% Minor insight (note.insight): small glyph, visibly indented
\newcommand{\noteheaderornament}[1]{%
  \par\nointerlineskip
  \vspace{0.12em}
  \begingroup
    % Tighten the row height so the ornament hugs the top padding
    \setlength{\baselineskip}{0.95\baselineskip}%
    \noindent\hspace{1.05em}\smash{\small #1}%
  \par\nointerlineskip
  \endgroup
}


% 3a. Badge strip header with neutral note color.
% Major insight (note.insight.major): badge strip aligned to left text edge
\newcommand{\noteheaderbadgenormal}[1]{%
  \begingroup
  \par\nointerlineskip
  \vspace{-0.25em}% counteract the restored top padding so the badge hugs the frame
  \noindent
  \expandafter\colorbox\expandafter{\notebadgebgnormal}{%
    \rule{0pt}{.10\baselineskip}%
    \hspace{0.35em}\raisebox{0.5ex}{\large #1}\hspace{0.35em}%
  }%
  \par\nointerlineskip
  \endgroup
}

% 3c. Badge strip header with breadcrumb accent color.
\newcommand{\noteheaderbadgebreadcrumb}[1]{%
  \begingroup
  \par\nointerlineskip
  \vspace{-0.25em}
  \noindent
  \expandafter\colorbox\expandafter{\notebadgebgbreadcrumb}{%
    \rule{0pt}{.10\baselineskip}%
    \hspace{0.35em}\raisebox{0.5ex}{\large #1}\hspace{0.35em}%
  }%
  \par\nointerlineskip
  \endgroup
}

% 3d. Badge strip header with navigation accent color.
\newcommand{\noteheaderbadgenavigation}[1]{%
  \begingroup
  \par\nointerlineskip
  \vspace{-0.25em}
  \noindent
  \expandafter\colorbox\expandafter{\notebadgebgnavigation}{%
    \rule{0pt}{.10\baselineskip}%
    \hspace{0.35em}\raisebox{0.5ex}{\large #1}\hspace{0.35em}%
  }%
  \par\nointerlineskip
  \endgroup
}

% 3e. Badge strip header with bridge accent color.
\newcommand{\noteheaderbadgebridge}[1]{%
  \begingroup
  \par\nointerlineskip
  \vspace{-0.25em}
  \noindent
  \expandafter\colorbox\expandafter{\notebadgebgbridge}{%
    \rule{0pt}{.10\baselineskip}%
    \hspace{0.35em}\raisebox{0.5ex}{\large #1}\hspace{0.35em}%
  }%
  \par\nointerlineskip
  \endgroup
}

% 3f. Badge strip header with diagnostic accent color.
\newcommand{\noteheaderbadgediagnostic}[1]{%
  \begingroup
  \par\nointerlineskip
  \vspace{-0.25em}
  \noindent
  \expandafter\colorbox\expandafter{\notebadgebgdiagnostic}{%
    \rule{0pt}{.10\baselineskip}%
    \hspace{0.35em}\raisebox{0.5ex}{\large #1}\hspace{0.35em}%
  }%
  \par\nointerlineskip
  \endgroup
}


% 3b. Badge strip header with warning note color.
\newcommand{\noteheaderbadgewarning}[1]{%
  \begingroup
  \par\nointerlineskip
  \vspace{-0.25em}
  \noindent
  \expandafter\colorbox\expandafter{\notebadgebgwarning}{%
    \strut\hspace{0.25em}\raisebox{0.05em}{#1}\hspace{0.25em}%
  }%
  \par\nointerlineskip
  \endgroup
}

% 3g. Badge strip header with misuse accent color (harsher warning tone).
\newcommand{\noteheaderbadgemisuse}[1]{%
  \begingroup
  \par\nointerlineskip
  \vspace{-0.25em}
  \noindent
  \expandafter\colorbox\expandafter{\notebadgebgmisuse}{%
    \strut\hspace{0.25em}\raisebox{0.05em}{#1}\hspace{0.25em}%
  }%
  \par\nointerlineskip
  \endgroup
}

% 3h. Badge strip header with closing accent (muted gray).
\newcommand{\noteheaderbadgeclosing}[1]{%
  \begingroup
  \par\nointerlineskip
  \vspace{-0.25em}
  \noindent
  \expandafter\colorbox\expandafter{\notebadgebgclosing}{%
    \rule{0pt}{.10\baselineskip}%
    \hspace{0.35em}\raisebox{0.5ex}{\large #1}\hspace{0.35em}%
  }%
  \par\nointerlineskip
  \endgroup
}


%--------------------------------------------------------------------
% Semantic header macros (map semantic types → presentation)
%--------------------------------------------------------------------

% note.idea  → inline lightbulb
\newcommand{\headernoteidea}{%
  \noteheaderinline{\glyphidea}%
}

% note.insight (minor) → ornamental asterism row
\newcommand{\headernoteinsight}{%
  \noteheaderornament{\glyphinsight}%
}

% note.insight.major (major) → badge strip with asterism
\newcommand{\headernoteinsightmajor}{%
  \noteheaderbadgenormal{\glyphinsight}%
}

% note.warning → badge strip with warning glyph and warning color
\newcommand{\headernotewarning}{%
  \noteheaderbadgewarning{\glyphwarning}%
}

% Optional: neutral note badge (for plain notes if needed)
\newcommand{\headernotenote}{%
  \noteheaderbadgenormal{\glyphnote}%
}

% note.breadcrumb → directional badge row
\newcommand{\headernotebreadcrumb}{%
  \noteheaderbadgebreadcrumb{\glyphbreadcrumb}%
}

% note.contrast → inline bidirectional arrow
\newcommand{\headernotecontrast}{%
  \noteheaderinline{\glyphcontrast}%
}

% note.navigation → teal navigation badge
\newcommand{\headernotenavigation}{%
  \noteheaderbadgenavigation{\glyphnavigation}%
}

% note.bridge → lavender badge emphasizing linkage
\newcommand{\headernotebridge}{%
  \noteheaderbadgebridge{\glyphbridge}%
}

% note.misuse → assertive red badge
\newcommand{\headernotemisuse}{%
  \noteheaderbadgemisuse{\glyphmisuse}%
}

% note.diagnostic → magnifying-glass badge
\newcommand{\headernotediagnostic}{%
  \noteheaderbadgediagnostic{\glyphdiagnostic}%
}

% note.closing → muted closing badge
\newcommand{\headernoteclosing}{%
  \noteheaderbadgeclosing{\glyphclosing}%
}


%--------------------------------------------------------------------
% Backwards-compatible macros expected by iasp-semantics.lua
%--------------------------------------------------------------------

% Inline lightbulb for note.idea
\newcommand{\notevariantideainline}{%
  \headernoteidea
}

% Minor insight header (note.insight)
\newcommand{\notevariantinsight}{%
  \headernoteinsight
}

% Major insight header (note.insight.major)
\newcommand{\notevariantinsightmajor}{%
  \headernoteinsightmajor
}

% Generic badge for notes where Lua inserts \notevariant{<variant>}:
% - "warning" → warning badge
% - anything else → neutral note badge
\newcommand{\notevariant}[1]{%
  \ifstrequal{#1}{warning}{%
    \headernotewarning
  }{%
    \ifstrequal{#1}{breadcrumb}{%
      \headernotebreadcrumb
    }{%
      \ifstrequal{#1}{contrast}{%
        \headernotecontrast
      }{%
        \ifstrequal{#1}{navigation}{%
          \headernotenavigation
        }{%
          \ifstrequal{#1}{bridge}{%
            \headernotebridge
          }{%
            \ifstrequal{#1}{misuse}{%
              \headernotemisuse
            }{%
              \ifstrequal{#1}{diagnostic}{%
                \headernotediagnostic
              }{%
                \ifstrequal{#1}{closing}{%
                  \headernoteclosing
                }{%
                  \headernotenote
                }%
              }%
            }%
          }%
        }%
      }%
    }%
  }%
}


%--------------------------------------------------------------------
% Note box environments
%--------------------------------------------------------------------

\newmdenv[
  skipabove=0.8\baselineskip,
  skipbelow=0.8\baselineskip,
  linewidth=0.8pt,
  linecolor=noteboxplainline,
  roundcorner=4pt,
  backgroundcolor=noteboxplainbg,
  innertopmargin=0.7em,
  innerbottommargin=0.7em,
  innerleftmargin=0.95em,
  innerrightmargin=0.95em,
]{noteboxplain}

\newmdenv[
  skipabove=0.8\baselineskip,
  skipbelow=0.8\baselineskip,
  linewidth=0.8pt,
  linecolor=noteboxline,
  roundcorner=4pt,
  backgroundcolor=noteboxbg,
  innertopmargin=0.7em,
  innerbottommargin=0.7em,
  innerleftmargin=0.95em,
  innerrightmargin=0.95em,
]{notebox}

\newmdenv[
  skipabove=0.8\baselineskip,
  skipbelow=0.8\baselineskip,
  linewidth=0.8pt,
  linecolor=noteboxbreadcrumbline,
  roundcorner=4pt,
  backgroundcolor=noteboxbreadcrumbbg,
  innertopmargin=0.7em,
  innerbottommargin=0.7em,
  innerleftmargin=0.95em,
  innerrightmargin=0.95em,
]{noteboxbreadcrumb}

\newmdenv[
  skipabove=0.8\baselineskip,
  skipbelow=0.8\baselineskip,
  linewidth=0.8pt,
  linecolor=noteboxwarningline,
  roundcorner=4pt,
  backgroundcolor=noteboxwarningbg,
  innertopmargin=0.7em,
  innerbottommargin=0.7em,
  innerleftmargin=0.95em,
  innerrightmargin=0.95em,
]{noteboxwarning}

\newmdenv[
  skipabove=0.8\baselineskip,
  skipbelow=0.8\baselineskip,
  linewidth=0.8pt,
  linecolor=noteboxcontrastline,
  roundcorner=4pt,
  backgroundcolor=noteboxcontrastbg,
  innertopmargin=0.7em,
  innerbottommargin=0.7em,
  innerleftmargin=0.95em,
  innerrightmargin=0.95em,
]{noteboxcontrast}

\newmdenv[
  skipabove=0.8\baselineskip,
  skipbelow=0.8\baselineskip,
  linewidth=0.8pt,
  linecolor=noteboxnavigationline,
  roundcorner=4pt,
  backgroundcolor=noteboxnavigationbg,
  innertopmargin=0.7em,
  innerbottommargin=0.7em,
  innerleftmargin=0.95em,
  innerrightmargin=0.95em,
]{noteboxnavigation}

\newmdenv[
  skipabove=0.8\baselineskip,
  skipbelow=0.8\baselineskip,
  linewidth=0.8pt,
  linecolor=noteboxbridgeline,
  roundcorner=4pt,
  backgroundcolor=noteboxbridgebg,
  innertopmargin=0.7em,
  innerbottommargin=0.7em,
  innerleftmargin=0.95em,
  innerrightmargin=0.95em,
]{noteboxbridge}

\newmdenv[
  skipabove=0.8\baselineskip,
  skipbelow=0.8\baselineskip,
  linewidth=0.8pt,
  linecolor=noteboxmisuseline,
  roundcorner=4pt,
  backgroundcolor=noteboxmisusebg,
  innertopmargin=0.7em,
  innerbottommargin=0.7em,
  innerleftmargin=0.95em,
  innerrightmargin=0.95em,
]{noteboxmisuse}

\newmdenv[
  skipabove=0.8\baselineskip,
  skipbelow=0.8\baselineskip,
  linewidth=0.8pt,
  linecolor=noteboxdiagnosticline,
  roundcorner=4pt,
  backgroundcolor=noteboxdiagnosticbg,
  innertopmargin=0.7em,
  innerbottommargin=0.7em,
  innerleftmargin=0.95em,
  innerrightmargin=0.95em,
]{noteboxdiagnostic}

\newmdenv[
  skipabove=0.8\baselineskip,
  skipbelow=0.8\baselineskip,
  linewidth=0.8pt,
  linecolor=noteboxclosingline,
  roundcorner=4pt,
  backgroundcolor=noteboxclosingbg,
  innertopmargin=0.7em,
  innerbottommargin=0.7em,
  innerleftmargin=0.95em,
  innerrightmargin=0.95em,
]{noteboxclosing}
